% LIVE-READY STAGING ONLY, 28 August 2026.
% Release gate: do not integrate this example until the arbitrary-field
% nonreduced exact-live statement/proof bytes receive exact-live verification.
% The rendered example below deliberately contains no gate-status language.

\begin{example}
\label{ex:p21p3-reducible-node-bt}
Let \(K\) be a complete nontrivially valued non-Archimedean field and let
\[
                 X=\{XY=0\}\subset\mathbf P^2_K.
\]
This is the reduced union of the two projective lines
\(L_1=\{X=0\}\) and \(L_2=\{Y=0\}\), meeting transversely at the
\(K\)-rational node \(o=[0:0:1]\).  Its analytification is good, separated,
strict, pure of dimension one, and boundaryless.  The normalization is
\[
 \nu:L_1^{\mathrm{an}}\amalg L_2^{\mathrm{an}}
          \longrightarrow X^{\mathrm{an}},
\]
write \(\nu_i\) for its restriction to \(L_i^{\mathrm{an}}\).  The two
points \(o_i\in L_i^{\mathrm{an}}\) above \(o\) are rigid.

By \cref{prop:p21p3-curve-rigidity}, every positive closed current of
bidegree \((0,0)\) on \(X^{\mathrm{an}}\) is
\begin{equation}
                         T=c_1\delta_{L_1}+c_2\delta_{L_2},
                         \qquad c_1,c_2\geq0.
 \label{eq:p21p3-node-current}
\end{equation}
The coefficients need not agree: the node does not create a balancing law
between distinct irreducible components.

Let \(W\subseteq X^{\mathrm{an}}\) be open, let \(u,v\) be continuous psh
functions on \(W\), and assume
\(
 \mathrm d'\mathrm d''[u]=\mathrm d'\mathrm d''[v].
\)
Write \(W_i=\nu_i^{-1}(W)\) and
\[
 \mu_i^u=\mathrm d'\mathrm d''[\nu_i^*u],
 \qquad
 \mu_i^v=\mathrm d'\mathrm d''[\nu_i^*v].
\]
If \(o\notin W\), the normalization-sheet images are already disjoint over
\(W\), and the curvature equality is sheetwise.  Suppose that \(o\in W\).
Away from \(o\), the images of the two normalization sheets are disjoint.
The curvature equality therefore shows that the signed measure
\(\mu_i^u-\mu_i^v\) is supported at \(o_i\) for each \(i\).  By
\cref{lem:p21p3-gfc-rigid-zero-atom}, neither positive measure
\(\mu_i^u\) nor \(\mu_i^v\) has an atom at the rigid point \(o_i\).  Hence
\begin{equation}
                              \mu_i^u=\mu_i^v
                              \qquad(i=1,2).
 \label{eq:p21p3-node-sheetwise-equality}
\end{equation}
Using \eqref{eq:p21p3-node-current}, proper pushforward from the two
normalization sheets, and \eqref{eq:p21p3-node-sheetwise-equality}, one
obtains
\[
 \mathrm d'\mathrm d''\bigl((u-v)T|_W\bigr)
   =\sum_{i=1}^2c_i(\nu_i)_*(\mu_i^u-\mu_i^v)=0.
\]
Thus equality of the two curvature measures cannot arise from cancellation
of opposite atoms on the two branches of the node.

There is also a concrete nonreduced variant.  Put
\[
 \mathfrak X^{(2)}
   :=\operatorname{Proj}K[X,Y,Z]/\bigl((XY)^2\bigr).
\]
Its reduction is \(X\), and its generic multiplicity along each \(L_i\) is
two.  By \cref{prop:p21p3-curve-rigidity-nonreduced}, its analytification
has, under the common current identification, the same ordinary CLD forms,
currents, positive cone, and continuous-pair Bedford--Taylor predicate as
\(X^{\mathrm{an}}\), whereas its fundamental integration current is
\begin{equation}
        \delta_{(\mathfrak X^{(2)})^{\mathrm{an}}}
          =2\delta_{L_1}+2\delta_{L_2}.
 \label{eq:p21p3-double-node-fundamental-current}
\end{equation}
Thus nilpotents change the generic cycle multiplicities, not the normalized
component currents that occur in \eqref{eq:p21p3-node-current}.

The example also shows why continuity and the rigid-point zero-atom input
are used in the descent: the argument establishes the continuous-pair
Bedford--Taylor property, not the literal possibly discontinuous scope of
\cite[Definition~8.6.2]{CLD25}.
\end{example}
